\documentclass[11pt]{amsart}
\usepackage[margin=1in]{geometry}
\usepackage{amssymb,mathtools}
\usepackage[hidelinks]{hyperref}
\hypersetup{pdftitle={Landau--Siegel zeros and a Stepanov-type nonvanishing lemma},
  pdfauthor={GPT-6 Pro (OpenAI) and Claude Opus 5.5 (Anthropic), posted and maintained by Jacob Pescador},
  pdfsubject={Preprint, version 2, 21 October 2026. AI-produced, not refereed.}}

\newtheorem{theorem}{Theorem}
\newtheorem{corollary}[theorem]{Corollary}
\newtheorem{lemma}[theorem]{Lemma}
\newtheorem{proposition}[theorem]{Proposition}
\theoremstyle{remark}
\newtheorem{remark}[theorem]{Remark}
\emergencystretch=1em
\raggedbottom
\makeatletter\@secpenalty=0\makeatother
\DeclareTextSymbolDefault{\textsection}{OMS}

\newcommand{\Z}{\mathbb{Z}}
\newcommand{\Q}{\mathbb{Q}}
\newcommand{\F}{\mathbb{F}}
\newcommand{\cO}{\mathcal{O}}
\newcommand{\cM}{\mathcal{M}}
\newcommand{\cL}{\mathcal{L}}
\newcommand{\sd}{\sqrt{d}}
\newcommand{\C}{\mathbb{C}}
\DeclarePairedDelimiter\abs{\lvert}{\rvert}

\title[Landau--Siegel zeros and a Stepanov-type lemma]{Landau--Siegel zeros and a\\ Stepanov-type nonvanishing lemma}
\author[GPT-6 Pro and Claude Opus 5.5]{GPT-6 Pro (OpenAI) and Claude Opus 5.5 (Anthropic)}
\dedicatory{Posted and maintained by Jacob Pescador}
\thanks{Version~2, 21 October 2026. All versions and corrections:\newline\url{https://jacobpescador.github.io/ai-math/papers/landau-siegel-stepanov/}}

\def\UrlBigBreaks{}\def\UrlBreaks{\do\/}\mathchardef\UrlBreakPenalty=100
\begin{document}

\begin{abstract}
Let $\chi$ be a primitive real character of conductor $q\ge3$. We prove Barreto's bound
\[
\sum_{\substack{p\le X\\ \chi(p)=-1}}\frac{\log p}p\le\frac12\log X+\frac14\log q+\log\log X+O(1)
\]
uniformly in $X\ge3$ and $q$, with an effective implied constant, and hence that every zero
$\beta\in(0,1)$ of $L(s,\chi)$ satisfies $(1-\beta)\log q\ge1-O(\log\log q/\log q)$, so there is no
Landau--Siegel zero. Our proof uses a single congruence for both the divisibility and the nonvanishing
of an interpolation determinant. As in OpenAI's proof that such zeros do not exist, the norm of the
determinant is bounded above by the size of its entries and below by its divisibility by high powers
of the odd inert primes $p$. Here the points are conjugate pairs $(\alpha,\alpha')$ of quadratic
integers, and conjugation agrees with the $p$th power map modulo $p$. Modulo a large inert prime $p_0$, every point lies on the
graph $w=z^{p_0}$ of Frobenius, and as $z^{p_0}$ has derivative zero modulo $p_0$, a linear relation
among the rows would give a polynomial in $z$ with more roots, counted with multiplicity, than its
degree. This is the multiplicity count of Stepanov's method. The points form Barreto's rectangular box,
which gives the $\frac14$.
\end{abstract}

\maketitle

\section{Introduction}

A Landau--Siegel zero is a real zero $\beta$ of a quadratic Dirichlet $L$-function $L(s,\chi)$ with
$(1-\beta)\log q$ small, $q$ being the conductor. Other Dirichlet $L$-functions have none~\cite[Ch.~14]{Davenport}.
Excluding such zeros gives the effective bound
$L(1,\chi)\gg1/\log q$ in place of Siegel's ineffective $L(1,\chi)\gg_\varepsilon q^{-\varepsilon}$, and so,
when $\chi$ is odd, an effective bound $\gg\sqrt q/\log q$ for the class number of $\Q(\sqrt{-q})$. Such zeros
were recently excluded by OpenAI~\cite{OpenAI}, as expounded by Banks~\cite{Banks}, and since then by Kerr~\cite{Kerr},
Green~\cite{Green}, Barreto~\cite{Barreto} and van Dijk~\cite{vanDijk}. All of these proofs rest on
divisibility by high powers of the inert primes, and they differ in their points, weights and
nonvanishing arguments.
This paper proves nonvanishing with the congruence that gives the divisibility, by a Stepanov-type
root count at an auxiliary inert prime, and lets the weight grow with $X$, which makes the bound uniform
in $X$ and $q$. With Barreto's rectangular box of points, the argument gives his bounds.

Let $\chi$ be a primitive real Dirichlet character of conductor $q\ge3$, and let $S_\pm(X)$ denote
the sum of $(\log p)/p$ over the primes $p\le X$ with $\chi(p)=\pm1$. Implied constants are
absolute and effective. We use Mertens' theorem $\sum_{p\le x}(\log p)/p=\log x+O(1)$ and Chebyshev's
bound $\sum_{p\le x}\log p\ll x$.

\begin{theorem}[Barreto~\cite{Barreto}]\label{thm:main}
For $X\ge3$ we have
\[
S_-(X)\le\frac12\log X+\frac14\log q+\log\log X+O(1).
\]
\end{theorem}

\begin{corollary}[Barreto~\cite{Barreto}]\label{cor:siegel}
Every zero $\beta\in(0,1)$ of $L(s,\chi)$ satisfies
\[
(1-\beta)\log q\ge1-O\Bigl(\frac{\log\log q}{\log q}\Bigr)\qquad\textup{and}\qquad(1-\beta)\log q\gg1.
\]
\end{corollary}


A real zero close to $1$ would make $\chi(p)=-1$ for most primes in a long range, and so make
$S_-(X)$ nearly $\log X$, contrary to Theorem~\ref{thm:main}. For fixed $q$, Mertens' theorem and
$L(1,\chi)\ne0$ give $S_-(X)=\frac12\log X+O_q(1)$, so the $\frac12$ cannot be lowered,
and the content of Theorem~\ref{thm:main} is its uniformity in $q$.

\emph{The idea.} Write $\chi$ as the Kronecker symbol $(d/\cdot)$ with $d$ a fundamental discriminant, $\abs{d}=q$, and for
$\alpha=x+y\sd$ put $\alpha'=x-y\sd$. If $p$ is odd and $\chi(p)=-1$, so that $p$ is inert
in $\Q(\sd)$, then conjugation on $\Z[\sd]/p$ is the Frobenius, so at every point $(z,w)=(\alpha,\alpha')$
\begin{equation}\label{eq:fact}
w\equiv z^p\pmod p,\qquad\text{and}\qquad\partial_z^\kappa(z^p-w)\equiv0\pmod p\quad(\kappa\ge1).
\end{equation}

The plan is to build a determinant $\delta$ with entries of controlled size that is a nonzero
algebraic integer divisible by high powers of many inert primes. A nonzero integer has absolute value at least the product of the
prime powers dividing it, so the divisibility bounds the norm of $\delta$ from below, Hadamard's
inequality bounds it from above, and comparing the two gives an upper bound for $S_-$.

Fix an inert prime $p_0$ and a weight $\ell$, which keeps the $w$-degrees small (they add to the size of the
entries but not to the divisibility, which comes from the $z$-degrees), and form the matrix with a row for each monomial $z^iw^j$
with $i+\ell j\le p_0-1$, and a column for each point $(\alpha,\alpha')$, with $\alpha$ in the box
$0\le x<N_1$, $0\le y<N_2$, and each order of $z$-derivative up to the largest $w$-degree, about $p_0/\ell$, as the root
count below requires. Each entry is that
derivative of the monomial at the point. Take $\delta$ to be a minor that uses every row, and $\delta'$
its conjugate, so that $\delta\delta'$ is an integer. There are about $p_0^2/(2\ell)$ rows but only about $p_0$ points in the box, and the
derivatives supply the missing columns without enlarging the box of points.

\emph{Divisibility.} At an inert prime $p$ with $\ell<p<p_0$, the change of basis
$z^iw^j\mapsto z^v(z^p-w)^hw^j$, where $i=ph+v$ with $0\le v<p$, trades $z$-degree for $w$-degree and leaves $\delta$
unchanged. By~\eqref{eq:fact} the new row is divisible by $p^h$, so $\delta$ is divisible by
about $p^{I/p}$, where $I=\sum i$.

\emph{Nonvanishing.} At $p_0$ every $z$-degree is below $p_0$, so there is nothing to rewrite,
and~\eqref{eq:fact} is used a second time. Modulo $p_0$ the points lie on the graph $w=z^{p_0}$ of
Frobenius, and since $\partial_z(z^{p_0}-w)\equiv0$, differentiating in $z$ commutes with substituting
$w=z^{p_0}$. If the box holds $p_0$ points distinct modulo $p_0$, a relation among the rows modulo $p_0$ therefore
becomes a polynomial in $z$ with more roots,
counted with multiplicity, than its degree (Lemma~\ref{lem:stepanov}), so $\delta$ can be chosen nonzero.
The derivatives thus serve both steps: they keep the divisibility at the small inert primes, and at $p_0$ they
supply the multiplicities. This argument uses characteristic $p_0$. Over $\C$, the
substitution $w=g(z)$ commutes with $\partial_z$ only for constant $g$, and each line $w=c$ contains at
most one of the points. Modulo $p_0$ the choice $g=z^{p_0}$ is also available, and its graph contains
all of them.

\emph{Size.} Call $\max(\abs{\alpha},\abs{\alpha'})$ the height of the point $(\alpha,\alpha')$. It
controls the entries of $\delta$ and $\delta'$, and the root count requires $p_0$ points, so we choose the
box to minimise their height. A box with sides $N_1,N_2$ has $N_1N_2$ points, each of height at most
$N_1+N_2\sqrt q$. The best shape balances the two terms: $N_1\approx p_0^{1/2}q^{1/4}$ and
$N_2\approx p_0^{1/2}q^{-1/4}$, with height $H\approx p_0^{1/2}q^{1/4}$. The points stay distinct modulo
$p_0$ when $q\le p_0^2$. In that range,
up to constants, no set of $p_0$ points of $\Z[\sd]$ does better, since
$\#\{\alpha:\max(\abs{\alpha},\abs{\alpha'})\le H\}\ll(1+H)(1+H/\sqrt q)$.

Heuristically, Hadamard's inequality gives $\log\abs{\delta\delta'}\lesssim2I\log H$, and the
divisibility gives $\log\abs{\delta\delta'}\gtrsim2I\sum(\log p)/p$ over the inert $p\in(\ell,p_0)$, so
\[
\sum_{\substack{\ell<p<p_0\\ \chi(p)=-1}}\frac{\log p}p\lesssim\log H\approx\frac12\log p_0+\frac14\log q.
\]
The right side is at most $\log p_0$, the bound given by Mertens' theorem, exactly when $q\le p_0^2$.
Taking $p_0$ to be the largest inert prime up to $X$, Mertens' theorem alone gives Theorem~\ref{thm:main} for larger $q$.

In the proof, the $w$-degrees and the derivatives add $O(\log p_0/\ell)$ to this bound, and the primes
$p\le\ell$, to which the rewriting does not apply, add $\log\ell$. We take $\ell\approx\log p_0$. Without derivatives, about $p_0^2/(2\ell)$ points would be needed,
of height about $p_0q^{1/4}/\sqrt\ell$. Adding $\log\ell$, this would give only
$S_-\lesssim\log p_0+\frac14\log q+\frac12\log\ell$, worse than Mertens' bound. The $\frac14$ gives the
constant $1$ in Corollary~\ref{cor:siegel} (Remark~\ref{rem:const}).

\emph{Why inert primes.} At an odd split prime, Frobenius is the identity on $\Z[\sd]/p\cong\F_p\times\F_p$,
while conjugation swaps the two factors, so no polynomial $g$ gives $\alpha'\equiv g(\alpha)\pmod p$ for all $\alpha\in\Z[\sd]$. The
congruence $z^p\equiv z$ that does hold does not survive differentiation, since $z^p-z$ has $z$-derivative
$-1$. The argument therefore bounds $S_-$ alone, which is the sum that a Landau--Siegel zero makes large.


\section{Two lemmas}\label{sec:lemmas}

Let $\cO=\Z[\sd]=\Z\oplus\Z\sd$.

\begin{lemma}\label{lem:frob}
Let $p$ be an odd prime with $\chi(p)=-1$. Then $\cO/p\cO$ is a field, and
$\alpha^p\equiv\alpha'\pmod{p\cO}$ for every $\alpha\in\cO$.
\end{lemma}

\begin{proof}
As $(d/p)=\chi(p)=-1$, $d$ is not a square modulo $p$, so $\cO/p\cO\cong\F_p[T]/(T^2-d)$ is a field. By the binomial theorem, Fermat's little
theorem and Euler's criterion, $(x+y\sd)^p\equiv x+yd^{(p-1)/2}\sd\equiv x-y\sd\pmod{p\cO}$. (Oddness is needed: $\cO$ need not be
maximal at $2$.)
\end{proof}

\begin{lemma}[after Stepanov~\cite{Stepanov}]\label{lem:stepanov}
Let $k$ be a field of characteristic $p$, let $1\le t\le p$, and let $F\in k[z,w]$ have $\deg_zF<p$
and $\deg_wF<t$. If $(\partial_z^\kappa F)(a,a^p)=0$ for $\kappa<t$ at $p$ distinct $a\in k$, then $F=0$.
\end{lemma}

\begin{proof}
Suppose $F\ne0$ and put $G(z)=F(z,z^p)$. The monomial $z^iw^j$ becomes $z^{i+pj}$, and these
exponents are distinct because $0\le i<p$, so $G\ne0$ and $\deg G\le(p-1)+p(t-1)=pt-1$. As the
derivative of $z^p$ vanishes in characteristic $p$, the chain rule gives
$G^{(\kappa)}(z)=(\partial_z^\kappa F)(z,z^p)$, so $G^{(\kappa)}(a)=0$ for $\kappa<t$ at each of the
$p$ points. Writing $G=\sum_\kappa c_\kappa(z-a)^\kappa$ gives $G^{(\kappa)}(a)=\kappa!\,c_\kappa$, and
$\kappa!\ne0$ in $k$ for $\kappa<t\le p$, so $(z-a)^t$ divides $G$ for each of the $p$ points. The $(z-a)^t$ are pairwise coprime, contradicting
$\deg G<pt$.
\end{proof}

Lemma~\ref{lem:stepanov} is the multiplicity count of Stepanov's method. Stepanov
constructs a nonzero auxiliary polynomial vanishing to high order at a set of points and uses the count to bound
the size of the set. Here the points are given, and the count shows that $F=0$. The count is exact on the rectangle $i<p$, $j<t$: the exponents
$i+pj$ fill $0,\dots,pt-1$ once each, so by Hermite interpolation the $pt$ values
$(\partial_z^\kappa F)(a,a^p)$ can be arbitrary. In \S\ref{sec:main} the lemma is applied
with $p=p_0$ to the triangle $i+\ell j\le D$ instead, which, unlike the rectangle, is preserved by the change
of basis at each inert prime in $(\ell,D]$.

\section{Proof of Theorem~\ref{thm:main}}\label{sec:main}

\begin{proposition}\label{prop:ell}
Let $p_0\ge5$ be a prime with $\chi(p_0)=-1$ and $q\le p_0^2$, put $D=p_0-1$, and let $\ell\ge2$ be
an integer. Then
\[
S_-(D)\le\Bigl(\frac12+\frac{7}{2\ell}\Bigr)\log D+\frac{\ell+1}{4\ell}\log q+\log\ell+O(1).
\]
\end{proposition}

\begin{proof}[Deduction of Theorem~\ref{thm:main}]
Let $p_0$ be the largest prime $\le X$ with $\chi(p_0)=-1$. If there is no such prime, or $p_0\le3$,
then $S_-(X)=O(1)$. Otherwise $S_-(X)=S_-(p_0)$ by maximality. If $q>p_0^2$, then Mertens' theorem gives
$S_-(X)\le\log p_0+O(1)\le\frac12\log X+\frac14\log q+O(1)$. If $q\le p_0^2$, then $S_-(X)=S_-(D)+O(1)$, and we apply
Proposition~\ref{prop:ell} with $\ell=\lceil\log p_0\rceil\ge2$. As $\log D\le\ell$ and
$\log q\le2\ell$, the terms $7\log D/(2\ell)$ and $\log q/(4\ell)$ are $O(1)$, and
$\log\ell\le\log\log X+1$.
\end{proof}

\begin{proof}[Proof of Proposition~\ref{prop:ell}]

\emph{Monomials.} Let $s=\lfloor D/\ell\rfloor$, and let $\cM$ be the set of $(i,j)\in\Z_{\ge0}^2$
with $i+\ell j\le D$, so that $i\le D<p_0$ and $j\le s$. Put $m=\abs{\cM}$, $I=\sum_\cM i$ and
$J=\sum_\cM j$. Reflecting $i\mapsto D-\ell j-i$ for each fixed $j$ gives $2I+\ell J=Dm$. Writing
$i=\ell u+r$ with $0\le r<\ell$, for each $r$ the pairs $(u,j)$ fill a triangle
$u+j\le\lfloor(D-r)/\ell\rfloor$, symmetric in $u$ and $j$, so $\sum u=\sum j$ and
$I=\sum(\ell u+r)\ge\ell\sum u=\ell J$. Therefore, as
$Dm=2I+\ell J\le3I$,
\begin{equation}\label{eq:counts}
J\le I/\ell,\qquad m\le3I/D.
\end{equation}

\emph{A nonzero algebraic integer.} Let $N_1=\lceil p_0^{1/2}q^{1/4}\rceil$ and
$N_2=\lceil p_0^{1/2}q^{-1/4}\rceil$. Then $N_1N_2\ge p_0$. Moreover, $N_1\le p_0$, because
$p_0^{1/2}q^{1/4}\le p_0$ when $q\le p_0^2$, and $N_2\le\lceil p_0^{1/2}\rceil\le p_0$. Let $A$
be the matrix over $\cO$ with rows indexed by $\cM$, ordered by increasing $i$, columns indexed by
$(x,y,\kappa)$ with $0\le x<N_1$, $0\le y<N_2$ and $0\le\kappa\le s$, and entries
\[
A_{(i,j),(x,y,\kappa)}=\partial_z^\kappa\bigl(z^iw^j\bigr)(\alpha,\alpha')
=(i)_\kappa\,\alpha^{i-\kappa}\alpha'^{j}\quad(\kappa\le i),\qquad\alpha=x+y\sd,
\]
and $0$ if $\kappa>i$, where $(i)_\kappa=i(i-1)\cdots(i-\kappa+1)$. Its rows are linearly
independent modulo $p_0$. Indeed, a relation among them over the field $\cO/p_0\cO$ is a polynomial
$F=\sum_\cM c_{ij}z^iw^j$ with $(\partial_z^\kappa F)(\bar\alpha,\bar\alpha^{p_0})=0$ for
$\kappa\le s$ (Lemma~\ref{lem:frob}) at the $N_1N_2\ge p_0$ reductions $\bar\alpha$, which are
distinct because $N_1,N_2\le p_0$ and $1,\sd$ are linearly independent over $\F_{p_0}$. As
$\deg_zF\le D<p_0$, $\deg_wF\le s$ and $s+1\le p_0$, Lemma~\ref{lem:stepanov} with $t=s+1$ gives
$F=0$. Hence we may fix an $m\times m$ minor $\delta\notin p_0\cO$ of $A$, which does not depend on the
primes $p$ below. Let $\delta'$ be its image under $\sd\mapsto-\sd$, which is the same minor with
$\alpha$ and $\alpha'$ interchanged. Then $n=\abs{\delta\delta'}$ is a positive integer, as
$\delta\delta'\in\Z\setminus\{0\}$. (We pass to the norm because $\abs{\delta}$ itself can be small when $d>0$.)

\emph{Its size.} Put $H=N_1+N_2\sqrt q$. As $\abs{\alpha},\abs{\alpha'}\le x+y\sqrt q<H$ and
$(i)_\kappa\le D^s$, every entry in row $(i,j)$ of the matrices of $\delta$ and $\delta'$ has modulus
at most $D^sH^{i+j}$, so that row has length at most $\sqrt m\,D^sH^{i+j}$. Hadamard's inequality, which
bounds a determinant by the product of its row lengths, gives
\begin{equation}\label{eq:upper}
\log n\le2(I+J)\log H+2ms\log D+m\log m.
\end{equation}

\emph{Its divisibility.} Let $p$ be a prime with $\ell<p\le D$ and $\chi(p)=-1$. As $\ell\ge2$, $p$ is
odd, so Lemma~\ref{lem:frob} applies. For $(i,j)\in\cM$ write $i=ph+v$ with $0\le v<p$, and replace
the row of $z^iw^j$ by that of
\[
z^v(z^p-w)^hw^j=z^iw^j+\sum_{\nu=1}^{h}(-1)^\nu\binom h\nu z^{i-p\nu}w^{j+\nu}.
\]
The monomials in the sum have smaller $z$-degree and smaller $i+\ell j$, namely $i+\ell j-(p-\ell)\nu$, so they
lie in $\cM$. Doing this for every
row adds to each row integer multiples of rows of smaller $i$, a unitriangular change, so $\delta$,
which uses every row, is unchanged. By Leibniz's rule,
$\partial_z^\kappa\bigl(z^v(z^p-w)^hw^j\bigr)$ is an integer combination of products each containing
$h$ factors $\partial_z^\lambda(z^p-w)$ with $\lambda\ge0$. Each such factor lies in $p\cO$ at
$(\alpha,\alpha')$. For $\lambda=0$ this is Lemma~\ref{lem:frob}. For $1\le\lambda\le p$ it holds because
$\partial_z^\lambda z^p=(p)_\lambda z^{p-\lambda}$ and $p\mid(p)_\lambda$, and for $\lambda>p$ the factor is
zero. So the new row of $(i,j)$ lies in $p^h\cO$, hence $\delta\in p^{E_p}\cO$. Conjugating,
$\delta'\in p^{E_p}\cO$, so $p^{2E_p}\mid n$ with
\begin{equation}\label{eq:div}
E_p=\sum_{(i,j)\in\cM}\Bigl\lfloor\frac ip\Bigr\rfloor\ge\frac Ip-m.
\end{equation}

\emph{Comparison.} As the pairwise coprime $p^{2E_p}$ all divide $n$, multiplying~\eqref{eq:div}
by $2\log p$, summing over $p$ and using Chebyshev's bound $\sum_{p\le D}\log p\ll D$ give
\[
2I\sum_{\substack{\ell<p\le D\\ \chi(p)=-1}}\frac{\log p}p\le2\sum_pE_p\log p+O(mD)\le\log n+O(mD).
\]
Insert~\eqref{eq:upper} and divide by $2I$. By~\eqref{eq:counts}, we have
$(I+J)/I\le1+1/\ell$, $ms/I\le3/\ell$ and $mD/I\le3$, and also $(m\log m)/I\le3(\log m)/D=O(1)$ since
$m\le(D+1)^2$. Mertens'
theorem gives $\log\ell+O(1)$ for the primes $p\le\ell$, so
\begin{equation}\label{eq:shape}
S_-(D)\le\log H+\frac{\log H+3\log D}\ell+\log\ell+O(1).
\end{equation}
Finally, $N_1\le p_0^{1/2}q^{1/4}+1$ and
$N_2\sqrt q\le p_0^{1/2}q^{1/4}+\sqrt q$, and $1\le\sqrt q\le p_0^{1/2}q^{1/4}$ as $q\le p_0^2$.
Hence $H\le4p_0^{1/2}q^{1/4}$ and $\log H\le\frac12\log D+\frac14\log q+O(1)$, and
inserting this in~\eqref{eq:shape} gives
\[
S_-(D)\le\Bigl(1+\frac1\ell\Bigr)\Bigl(\frac12\log D+\frac14\log q\Bigr)+\frac3\ell\log D+\log\ell+O(1),
\]
which is the proposition, as $\frac12+\frac1{2\ell}+\frac3\ell=\frac12+\frac7{2\ell}$.
\end{proof}

\section{Proof of Corollary~\ref{cor:siegel}}\label{sec:cor}

A zero $\beta$ near $1$ makes $-L'/L(\sigma,\chi)$ very negative just to the right of $1$, while
Theorem~\ref{thm:main} bounds it below. We track the constants, as Barreto~\cite{Barreto} does.

\begin{proof}[Proof of Corollary~\ref{cor:siegel}]
\emph{Lower bound for $-L'/L$.} Let $T(x)=\sum_{p\le x}\chi(p)(\log p)/p=S_+(x)-S_-(x)$, so that $T(x)=0$ for
$1\le x<2$. The primes with $\chi(p)=0$ divide $q$, and
\[
\sum_{p\mid q}\frac{\log p}p\le\sum_{p\le\log q}\frac{\log p}p+\frac1{\log q}\sum_{p\mid q}\log p
\le\log\log q+O(1).
\]
So Mertens' theorem gives $T(x)\ge\log x-2S_-(x)-O(\log\log q)$. Writing $\log^+=\max(\log,0)$,
Theorem~\ref{thm:main} then gives some $K=\frac12\log q+O(\log\log q)$ with $K\ge\frac12\log q$ and
\begin{equation}\label{eq:Tlower}
T(x)\ge-K-2\log^+\log x\qquad(x\ge1),
\end{equation}
which is trivial for $x<3$, where $T(x)\ge-\frac12\log2>-K$. Let $0<u\le1$ and $\sigma=1+u$. The
logarithmic derivative is an average of $T$ over logarithmic scales. By partial summation, and since
the prime powers contribute $O(1)$ (as $\sum_p(\log p)/(p(p-1))$ converges),
\[
-\frac{L'}{L}(\sigma,\chi)=u\int_1^\infty T(x)\,x^{-u-1}\,dx+O(1).
\]
The weight $u\,x^{-u-1}$ has total mass $1$ and satisfies $u^2\int_1^\infty(\log x)\,x^{-u-1}\,dx=1$.
Since $\log y\le y$, we have $\log^+\log x\le\log(1/u)+u\log x$, and so~\eqref{eq:Tlower} gives
\begin{equation}\label{eq:Llower}
-\frac{L'}{L}(\sigma,\chi)\ge-\frac12\log q-2\log\frac1u-O(\log\log q).
\end{equation}

\emph{Upper bound for $-L'/L$.} Let $L(\beta,\chi)=0$ with $0<\beta<1$. Taking real parts in the formula for
$L'/L$ that comes from the Hadamard
product~\cite[Ch.~12, (17)--(18), and Ch.~14, p.~89, (2)]{Davenport},
\[
-\frac{L'}{L}(\sigma,\chi)=\frac12\log\frac q\pi
+\frac12\frac{\Gamma'}{\Gamma}\Bigl(\frac{\sigma+\mathfrak a}2\Bigr)
-\sum_\rho\operatorname{Re}\frac1{\sigma-\rho},\qquad \chi(-1)=(-1)^{\mathfrak a},\ \mathfrak a\in\{0,1\},
\]
where $\rho$ runs over the nontrivial zeros of $L(s,\chi)$, counted with multiplicity. The gamma
term is bounded, as $(\sigma+\mathfrak a)/2\in(\frac12,\frac32]$, and every term of the sum is
positive because $\operatorname{Re}\rho\le1<\sigma$. Keeping only $\rho=\beta$ and using~\eqref{eq:Llower},
\begin{equation}\label{eq:key}
\frac1{u+1-\beta}\le\log q+2\log\frac1u+O(\log\log q)\qquad(0<u\le1).
\end{equation}

\emph{Conclusion.} Put $\cL=\log q$ and $u=\cL^{-2}$, so that $u\cL=1/\cL$ and
$\log(1/u)=2\log\cL$. By~\eqref{eq:key},
\[
(1-\beta)\cL\ge\frac{\cL}{\cL+O(\log\cL)}-\frac1\cL\ge1-O\Bigl(\frac{\log\cL}{\cL}\Bigr).
\]
This bound is empty for small $q$. To get $(1-\beta)\log q\gg1$ for every $q$, take $u=\eta/\cL$
instead, with $0<\eta\le1$. Since $\log(1/u)=\log(1/\eta)+\log\cL\le\log(1/\eta)+\cL$,~\eqref{eq:key} gives
$1/(u+1-\beta)\ll(1+\log(1/\eta))\cL$. Hence $(1-\beta)\cL\ge c_0/(1+\log(1/\eta))-\eta$ with $c_0>0$ absolute, and
this is positive once $\eta$ is small enough.
\end{proof}

\begin{remark}\label{rem:const}
The constant $1$ in Corollary~\ref{cor:siegel} is $1/(\frac12+2\cdot\frac14)$: the $\frac12$ comes from the conductor in the Hadamard formula, and $2\cdot\frac14$ from Theorem~\ref{thm:main}, as $S_-$ appears twice in
$T=(S_++S_-)-2S_-$. The constant also depends on the coefficient of $\log X$ in Theorem~\ref{thm:main} being exactly
$\frac12$, so that $2S_-$ cancels the $\log x$ of Mertens' theorem. With $\frac12+\varepsilon$ in its
place, the right side of~\eqref{eq:key} would gain a term $2\varepsilon/u$, and the argument would give
a constant below $1$.
\end{remark}

\section{Relation to other proofs}\label{sec:rel}

The OpenAI proof~\cite{OpenAI}, expounded by Banks~\cite{Banks}, works over
$\Q(\sqrt{d_0},\sqrt2)$, where $d_0$ is the squarefree part of $d$. As both note, it is an
interpolation-determinant argument~\cite{Laurent}, with the smallness $p$-adic rather than
archimedean, as in Dobrowolski's work on Lehmer's problem~\cite{Dobrowolski} (a connection noted by
K.~Pratt in the comments to~\cite{Tao}), whose determinant form is due to Cantor and
Straus~\cite{CantorStraus}.

\emph{Divisibility.} All the proofs use the congruence $w\equiv z^p\pmod p$. Green~\cite{Green} works in
$\Q(\sd)$ with the points $(\alpha,\alpha')$, takes $z$-derivatives and applies Leibniz's rule, and orders the rows by
$z$-degree. We follow him in this. Where his main proof uses a Gram determinant, we use a nonzero minor,
as a remark in~\cite[version~2, \S2]{Green} also suggests. The
change of basis to $z^v(z^p-w)^h$ comes from~\cite{OpenAI}, in Kerr's form~\cite{Kerr}.

\emph{Points and weights.} A square box of points, as in Green's proof, gives $\frac12\log q$, and
Barreto's~\cite{Barreto} rectangular box gives $\frac14\log q$. A fixed weight, as in OpenAI's and Kerr's
proofs, gives the bound at $X$ a fixed power of $q$, and a weight growing with $X$ costs $\log\log X$.

\emph{Nonvanishing.} OpenAI~\cite{OpenAI} uses an interpolation lemma with a greedy choice of monomial rows.
Kerr~\cite{Kerr} constructs an auxiliary polynomial by the pigeonhole principle in place of a
determinant, and Masser's zero estimate~\cite{Masser} shows that it cannot vanish at all the points.
Green~\cite[Lemma~3.1]{Green} bounds multiplicities by resultants at any points with distinct
second coordinates satisfying an incidence bound, which Schwartz--Zippel supplies. Barreto~\cite{Barreto} uses a Wronskian and a minor in $\Z$. Van Dijk~\cite{vanDijk} proves nonvanishing by the same root count, and his version~2 makes
the argument explicit for class number one: an imaginary quadratic field of class number one has discriminant
of absolute value below $8\cdot10^9$.

\emph{Versions.} Version~1 of this paper~\cite{v1}, dated 18 October 2026, had the root count of Lemma~\ref{lem:stepanov} and the weight
growing with $X$, and with a square box it proved Theorem~\ref{thm:main} with $\frac12\log q$ in place of
$\frac14\log q$.

\section*{Authorship}
The mathematics and text were produced by the AI systems named above, building on their version~1~\cite{v1}
and the work cited above. Jacob Pescador initiated and coordinated the project and
maintains this posting.

\begin{thebibliography}{99}\footnotesize\raggedright\setlength{\itemsep}{0pt}
\bibitem{Banks} W.\,D. Banks, \emph{On the OpenAI proof that Landau--Siegel zeros do not exist},
preprint, version~1, 11 October 2026, \href{https://arxiv.org/abs/2610.15483v1}{arXiv:2610.15483v1}.
\bibitem{Barreto} K. Barreto, \emph{An elementary proof of the non-existence of Landau--Siegel
zeros}, posted on X (@AcerFur), 19 October 2026, unavailable as of 20 October 2026.
\bibitem{CantorStraus} D.\,C. Cantor and E.\,G. Straus, \emph{On a conjecture of D.\,H. Lehmer},
Acta Arith. \textbf{42} (1982), 97--100; correction, ibid. \textbf{42} (1983), 327.
\bibitem{Davenport} H. Davenport, \emph{Multiplicative Number Theory}, 2nd ed., revised by
H.\,L. Montgomery, GTM 74, Springer, 1980.
\bibitem{Dobrowolski} E. Dobrowolski, \emph{On a question of Lehmer and the number of irreducible
factors of a polynomial}, Acta Arith. \textbf{34} (1979), 391--401.
\bibitem{v1} GPT-6 Pro and Claude Opus 5.5, \emph{Landau--Siegel zeros and a Stepanov-type
nonvanishing lemma}, preprint, version~1, 18 October 2026,
\url{https://jacobpescador.github.io/ai-math/papers/landau-siegel-stepanov/v1/landau-siegel-stepanov.pdf}.
\bibitem{Green} B. Green, \emph{On Landau--Siegel zeros}, preprint, version~1, 17 October 2026,
\href{https://hexagonmath.org/abs/2610.00311v1}{hexagon:2610.00311v1}, and version~2, 18 October 2026,
\href{https://hexagonmath.org/abs/2610.00311v2}{hexagon:2610.00311v2}.
\bibitem{Kerr} B. Kerr, \emph{A simpler proof that there are no Landau--Siegel zeros}, preprint,
version~1, 15 October 2026, \href{https://arxiv.org/abs/2610.21099v1}{arXiv:2610.21099v1}.
\bibitem{Laurent} M. Laurent, \emph{Sur quelques r\'esultats r\'ecents de transcendance},
Ast\'erisque \textbf{198--199--200} (1991), 209--230.
\bibitem{Masser} D.\,W. Masser, \emph{On polynomials and exponential polynomials in several complex
variables}, Invent. Math. \textbf{63} (1981), 81--95.
\bibitem{OpenAI} OpenAI, \emph{Uniform exclusion of Landau--Siegel zeros}, preprint, 1 October 2026,
\href{https://github.com/openai/math/blob/adc7f1241b42e322a6451854ab7e4b4c146bf78a/preprints/Uniform-exclusion-of-Landau-Siegel-zeros-October-1-2026/paper.pdf}{OpenAI Math Release (GitHub)}.
\bibitem{Stepanov} S.\,A. Stepanov, \emph{On the number of points of a hyperelliptic curve over a
finite prime field}, Math. USSR-Izv. \textbf{3} (1969), 1103--1114.
\bibitem{Tao} T. Tao, \href{https://terrytao.wordpress.com/2026/10/19/a-sketch-of-the-proof-of-non-existence-of-siegel-zeroes/}{\emph{A sketch of the proof of non-existence of Siegel zeroes}}, What's new (blog), 19 October 2026; comment cited: K.~Pratt, \#comment-698411.
\bibitem{vanDijk} R. van Dijk, \emph{A Frobenius determinant proof of the non-existence of
Landau--Siegel zeros}, preprint, version~1, 20 October 2026,
\href{https://hexagonmath.org/abs/2610.00409v1}{hexagon:2610.00409v1}, and version~2, 21 October 2026,
\href{https://hexagonmath.org/abs/2610.00409v2}{hexagon:2610.00409v2}.
\end{thebibliography}

\end{document}
